\documentclass[a4paper,11pt]{article}
\usepackage{graphicx}
\usepackage{amsmath}
\usepackage{amssymb}
\usepackage{hyperref}
\usepackage{booktabs}
\usepackage{caption}
\usepackage{geometry}
\usepackage{authblk}  % Package for author and affiliation formatting

\geometry{a4paper, margin=25mm}



\title{Title of the Paper}
% Authors and Affiliations
\author[1]{
DO NOT ADD NAMES OR AFFILIATIONS HERE
}

\author[2]{
DO NOT ADD NAMES OR AFFILIATIONS HERE. Authors will get an opportunity 
to add these after paper review.
Do not add email addresses here for review copy.
}


\affil[1]{Department Name, University Name, City, Country, Email: author1@example.com}
\affil[2]{Department Name, University Name, City, Country, Email: author2@example.com}

\date{}  % This removes the date

\renewcommand\Authands{ and }  % Fixes "and" between author names\begin{document}


\maketitle
\begin{document}
% --- INSERT SICET LOGO & CONFERENCE DETAILS HERE ---




\section*{Abstract}
In this paper, the formatting requirements for the 4th SLIIT International Conference on Engineering and Technology (SICET25) are described. Please review this document to learn about the formatting of text, table captions, references, and the method to include the indexing information. The conference proceedings will be published in an electronic format. The full paper in MS Word file shall be written in compliance with these instructions. At a later stage, it will be converted into Portable Document Format (PDF).

\textbf{Keywords:} 4 to 8 keywords, separated by commas.

\section{Introduction}
It is expected that authors will submit carefully written and proofread material. Careful checking for spelling and grammatical errors should be performed. The number of pages of the paper should be between 6 and 12.

Papers should clearly describe the background of the subject, the authors' work, including the methods used, results, and concluding discussion on the importance of the work. Papers are to be prepared in English and SI units must be used. Technical terms should be explained unless they may be considered to be known to the conference community. 

\section{Paper Format}
The uniform appearance will assist the reader in reading papers in the proceedings. It is therefore suggested to authors to use the example of this file to construct their papers. This particular example uses an A4 format with 25 mm margins left, right, top, and bottom and Times New Roman font.

All text paragraphs should be single spaced, with the first line indented by 10 mm. Double spacing should only be used before and after headings and subheadings as shown in this example. Position and style of headings and subheadings should follow this example. No spaces should be placed between paragraphs.

\section{Header, Footer, Page Numbering}
Authors are asked to replace the "XXX" number (with the paper code that was assigned when the paper was accepted) on the header of the first page and on the footer of all pages in order to set a unique page number in the Proceedings.

\section{Fonts}
Papers should use 11-point Times New Roman font. The styles available are bold, italic, and underlined. It is recommended that text in figures should not be smaller than a 10-point font size.

\section{Tables and Figures}
Table captions and figure captions should be sufficient to explain the table or figure without needing to refer to the text. The following is an example for Table~\ref{tab:example}.

\begin{table}[h]
\centering
\caption{Title of the Table}
\label{tab:example}
\begin{tabular}{ccc}
\toprule
Heading & Heading & Heading \\
\midrule
Information & Information & Information \\
Information & Information & Information \\
Information & Information & Information \\
\bottomrule
\end{tabular}
\end{table}

Leave a single space before and after a table or figure. Tables and figures should be placed close to their first reference in the text. All figures and tables should be numbered with Arabic numerals. Table headings should be centered above the tables. Figure captions should be centered below the figures.

\begin{figure}[h]
\centering
\includegraphics[width=0.7\textwidth]{example-image}
\caption{Number of patents on nanotechnology with time}
\label{fig:example}
\end{figure}


\begin{figure}[h]
\centering
\includegraphics[width=0.7\textwidth]{figure_1.png}
\caption{Figure 1 caption}
\label{fig:figure_1}
\end{figure}


\section{Equations}
Each equation should be presented on a separate line from the text with a blank space above and below. Equations should be clear, and expressions used should be explained in the text. The equations should be numbered consecutively at the outer right margin, as shown below.

\begin{equation}
E = mc^2
\end{equation}

\section{References}
Proper references must be included throughout the text and the list of references must be provided in this section. Please use the APA style for in-line citations and the listing of references.

\begin{thebibliography}{9}
\bibitem{ref1} Author(s) of the reference. (Year). \textit{Title of the work}. Publisher. DOI or URL.
\bibitem{ref2} Another reference example. (Year). \textit{Title of the work}. Journal Name, \textbf{Volume}(Issue), Pages.
\end{thebibliography}

\end{document}
